\section{Overview}
This document describes the MQTT protocol design for the distributed real-time telemetry system developed as part of the EGH490 thesis project at the Queensland University of Technology. The system acquires electrical measurements from STM32-based data acquisition nodes deployed across the QUT microgrid and transports them to a centralised time-series historian via an MQTT message broker.

The design covers broker selection and configuration, topic hierarchy, payload format, quality of service strategy, client architecture, security posture, and historian integration. All decisions are scoped to a two-node validation testbed operating over a wired Ethernet LAN, with the architecture designed to generalise to an arbitrary number of nodes without structural change.

Version 2 of this document records the system as built. Where the implementation departed from the original design, the as-built behaviour is described, and design elements that were not implemented are listed as deferred in the relevant section.
